\clearpage\section*{Solutions EV3-T2}\label{EV3-T2-CORR}
\textbf{MCQ:} 1 D, 2 D, 3 B, 4 D, 5 C, 6 C, 7 A, 8 A, 9 B, 10 A, 11 B, 12 A. Each correct answer: 0.5 point; total 6.
\paragraph{Checks of the fluency questions.} Order of operations; sum of fractions; subtraction of exponents; multiplying coefficient 1.2; division by 0.8; product 1.25 times 0.8; solving an affine equation; vertical increment divided by horizontal; mean; angle divided by 360; complement to 1; substitution in brackets.
\subsection*{Exercise 1 -- Simulation and mean}
\begin{enumerate}
\item $E(X)=0+\dfrac{1}{2}+\dfrac{3}{2}=2$. \textit{Marking: 1 point. Result and basic justification expected; half a point possible for a correct but unfinished approach.}
\item $E(X^2)=1+9=10$, variance 6, standard deviation $\sqrt6$. \textit{Marking: 2 points. Method or justification: 1; calculation and conclusion: 1.}
\item $[2-\dfrac{2\sqrt 6}{5};2+\dfrac{2\sqrt 6}{5}]$. A frequency over a finite number of simulations fluctuates. \textit{Marking: 2 points. Method or justification: 1; calculation and conclusion: 1.}
\end{enumerate}
\subsection*{Exercise 2 -- Comparing two models}
\begin{enumerate}
\item $h\prime=2(e^{2x}-1)$. \textit{Marking: 1 point. Result and basic justification expected; half a point possible for a correct but unfinished approach.}
\item Decreasing up to 0 then increasing; minimum 0. So $f\geq g$ with equality only at 0. \textit{Marking: 2 points. Method or justification: 1; calculation and conclusion: 1.}
\item Common ratio $e^2$; sum $1+e^2+e^4+e^6=\dfrac{1-e^8}{1-e^2}$. \textit{Marking: 2 points. Method or justification: 1; calculation and conclusion: 1.}
\end{enumerate}
\subsection*{Exercise 3 -- A displaced circle}
\begin{enumerate}
\item Centre (1;-1), radius 2. \textit{Marking: 1 point. Result and basic justification expected; half a point possible for a correct but unfinished approach.}
\item A(-1;-1), B(3;-1). \textit{Marking: 1 point. Result and basic justification expected; half a point possible for a correct but unfinished approach.}
\item AB is a diameter, so the product is 0, endpoints included; identity $MI^2-4=0$. \textit{Marking: 2 points. Method or justification: 1; calculation and conclusion: 1.}
\end{enumerate}
\paragraph{Calibration.} Target duration: 20--25 min for the fluency questions, 80--85 min for the exercises, 10--15 min for proofreading. These durations are editorial estimates to be measured in class. Accept any correct method; do not penalise twice a numerical error that is propagated consistently. The proposed marking scheme is not a decision of the team to harmonise marking.
